We now show that the length of the shortest nontrivial simple closed
curve on the cusp torus is bounded throughout the cone deformation,
provided the tube radius about the singular locus stays large.  Recall
here that we are calling this shortest curve the meridian of the
torus.  The following lemma is the last piece needed to complete the
proof of Theorem \ref{thm:main-norm-height}.

\begin{lemma}
Suppose $M$ is a hyperbolic cone manifold with one cusp, and with
singular locus $\Sigma$.  Suppose the tube radius $R$ about the
singular locus is at least $\log(2/\sqrt{3})$.  Then there exists a
horoball neighborhood $U$ of the cusp such that the meridian of
$\partial U$ has length at least $1-e^{-2R}$.
  \label{lemma:meridian-bound}
\end{lemma}

\begin{proof}
Let $X = M-\Sigma$.  $X$ admits an incomplete hyperbolic metric.
Expand the tube $V = \cup V_i$ about the
singular locus $\Sigma$ until the tube hits itself.  Each $V_i$ has
the same tube radius $R$.  Next expand the horoball $U$ about
the cusp until it either becomes tangent to itself or to $V$.

We obtain a region $\Omega$ by continuing to expand $U$ as follows.
First, choose a particular lift $U_a$ of $U$ in the universal cover.
Because $U$ is embedded in $\HH^3$ by the developing map on the
universal cover $\widetilde{X}$ of $X$, we can assume the developing
map is one to one on $U_a$ and takes $\partial U_a$ to the plane $z=1$
in $\HH^3$.

Take $\Omega= U_a \cup S$, where $S$ is the set consisting of all
points $x$ in $\tild{X}$ such that the distance $\mbox{dist}(x,
\partial U_a) < \mbox{dist}(x, \partial U_b)$ for $U_b$ any other lift
of $U$, and also such that the distance $\mbox{dist}(\phi(x), \partial
U) < \mbox{dist}(\partial U,\Sigma)$, where $\phi$ is the covering
projection. Note that $\Omega$ is embedded in $\HH^3$ by the
developing map.  That is, the developing map is one to one on
$\Omega$.

Let $p$ be a point on $\partial U$ closest to $\Sigma$, and let $p_0$
and $p_1$ be lifts of $p$ on $\partial U_a$ such that the projection
of the path from $p_0$ to $p_1$ is shortest on $\partial U$.  Let
$a_0$ be the point in the universal cover on a lift of $\Sigma$ such
that the path from $p_0$ to $a_0$ has length equal to the minimum
distance from $p_0$ to $\Sigma$.  Similarly, let $a_1$ be the point on
a lift of $\Sigma$ such that the path from $p_1$ to $a_1$ has length
equal to the minimum distance from $p_1$ to $\Sigma$.  Note the
interiors of the paths between $p_0$ and $a_0$ and between $p_1$ and
$a_1$ are contained in $\Omega$. 

\begin{figure}
  \begin{center}
  \input{figures/area_quad.pstex_t}
  \end{center}
  \caption{The top edge of the quadrilateral maps to the shortest path
  on $\partial U$ under the covering projection. }
  \label{fig:area-quad}
\end{figure}

Then note the interior of the geodesic $\beta$ from $a_0$ to $a_1$ is
also in $\Omega$, along with all the points in the quadrilateral with
sides given by $\beta$ and the paths from $p_0$ to $p_1$, from $p_0$
to $a_0$, and from $p_1$ to $a_1$.  This is because, first, all points
$x$ in this quadrilateral (except the endpoints of $\beta$) have
$\mbox{dist}(\phi(x), \partial U) < \mbox{dist}(\partial U, \Sigma)$.
Also, if any $x$ were closer to some other lift $U_b$ of $U$, then
there would be points in the quadrilateral equidistant to two
different lifts of $U$.  All points equidistant to two different lifts
of $U$ lie on some totally geodesic planes of $\HH^3$.  Hence some
geodesic must intersect $\beta$ but not the paths from $p_0$ to $a_0$,
$p_1$ to $a_1$, or the path on the horosphere $\partial U_a$ from
$p_0$ to $p_1$.  This is impossible.

Now let $H$ be the length of the paths from $p_0$ to $a_0$ and from
$p_1$ to $a_1$.  Let $L$ be the length of $\beta$.  We know that
$H\geq R$ and $L\geq 2R$ since the distance between any two components
of the lift of $\Sigma$ is greater than or equal to $2R$.  We let $d$
be the length of the path on the horosphere $\partial U_a$ between
$p_0$ and $p_1$.  Since $\Omega$ is embedded by the developing map, we
can assume the quadrilateral is of the form shown in Figure
\ref{fig:area-quad}.

Then solving for $d$ in terms of $H$ and $L$, $d= 2e^{-H}\sinh(L/2)$.

Case 1:  $H=R$.  Then $d \geq 2e^{-R}\sinh(R) = 1-e^{-2R}$.

Case 2: $H>R$ Then $U$ must be tangent to itself, say at a point $q$,
and not necessarily tangent to the tube about $\Sigma$.  Then the
image of $\Omega$ under the developing map must contain some portion
of the region above the geodesic surface equidistant from a horizontal
plane (the lift of $\partial U_a$) and some other horosphere.

This geodesic surface is some portion of the (Euclidean) hemisphere
with radius $1$, tangent to the lift of $\partial U_a$ at a lift $q_0$
of $q$.  See Figure \ref{fig:d-long}.  The image of $\Omega$ under the
developing map contains a neighborhood of points on this hemisphere
about $q_0$.  If it were to contain all the points on the hemisphere
between the dotted lines in Figure \ref{fig:d-long} (that is, points
whose projection to $\CC$ have Euclidean distance at most $1/2$ from
the projection of $q_0$ to $\CC$), then $d\geq 1$.

Let $x$ be one of these points.  That is, the projection of $x$ to
$\CC$ has Euclidean distance at most $1/2$ from the projection of
$q_0$ to $\CC$.  Note $x$ will lie in $\Omega$ provided
$\mbox{dist}(\phi(x), \partial U) < \mbox{dist}(\partial U, \Sigma) =
H$, or provided $H$ is large enough.  But now, $\mbox{dist}(\phi(x),
\partial U) = \log( 1/\mbox{height}(x))$, where $\mbox{height}(x)$
denotes the $z$--coordinate of $x$ in $\HH^3$.  Thus this distance is
at most $\log(2/\sqrt{3})$.

\begin{figure}
\input{figures/d-long.pstex_t}	
\caption{}
\label{fig:d-long}
\end{figure}

Since $H>R$, we can ensure $d\geq 1$ by restricting to $R\geq
\log(2/\sqrt{3})$.
\end{proof}

When we plug the result of Lemma \ref{lemma:meridian-bound} into
Theorem \ref{thm:norm-height-mer}, along with the bound on $\sum b_j$
of the previous section, we may complete the proof of Theorem
\ref{thm:main-norm-height}.

\medskip

\noindent\textbf{Theorem \ref{thm:main-norm-height}.} \emph{
  Let $M=X_\tau$ be a hyperbolic 3--manifold with a cusp, and let
  $X=X_0$ be a hyperbolic 3--manifold which can be joined to $M$ by a
  smooth family of hyperbolic cone manifolds $X_t$ with tube radius at
  least $R_1 \geq 0.56$, and with $n$ components of the singular
  locus. Let $U_t$ be a horoball neighborhood about
  the cusp.  Let $h(M)$ denote the normalized height of the cusp torus
  $\partial U_\tau$.  Let $h(X)$ denote the normalized height of
  $\partial U_0$.  Then the
%% Aug 6, 2007: parameterization:
	hyperbolic cone deformation can be parameterized such that the
%%
	change in normalized height is bounded in
  terms of $R_1$ alone:
$$-(2\pi)^2\frac{\sqrt{n\,C(R_1)}}{(1-e^{-2R_1})\sqrt{2}} \leq
 h(M) - h(X) \leq
(2\pi)^2\frac{\sqrt{n\,C(R_1)}}{(1-e^{-2R_1})\sqrt{2}}.$$
}

\begin{proof}
  Theorem \ref{thm:norm-height-mer}, combined with Theorem
  \ref{thm:bound-bdry-time} and Lemma \ref{lemma:meridian-bound}
  imply that
$$
  -\frac{\sqrt{n\,C(R_1)}}{(1-e^{-2R_1})\sqrt{2}} \leq \dot{h} \leq
\frac{\sqrt{n\,C(R_1)}}{(1-e^{-2R_1})\sqrt{2}}
$$
and that the total time of the deformation is $t=(2\pi)^2$.  We
integrate the above inequality over the deformation.  The bounds on
the left and right are independent of $t$.  Thus we obtain the
conclusion of the theorem.
\end{proof}
